\section{Completing the irrationality contradiction}\label{sec:contradiction}

We now prove Proposition~\ref{prop:criterion}. Assume its two
inequalities hold and suppose that at most $B-1$ of
$\beta(2),\ldots,\beta(s)$ are irrational. Choose $B-1$ distinct
even indices $i_1<\cdots<i_{B-1}$ containing all possible
irrational indices, adding further indices if necessary. There
are enough choices, since $s\ge DL\ge10D\ge10B$.

The coefficients $c_j$ have already fixed the common multiplier
$W_n$. We now choose integer weights $w=(w_1,\ldots,w_B)$ on the
$B$ forms to annihilate the columns
\[
 (1^{i_j},2^{i_j},\ldots,B^{i_j})^{\mathsf T},
                  \qquad j=1,\ldots,B-1,
\]
by using their signed maximal minors. Its last entry $w_B$ is nonzero.
Indeed, the corresponding minor is a generalized Vandermonde
determinant at the distinct positive points
$1^2,\ldots,(B-1)^2$, with exponents $i_j/2$. A singular minor
would give a nonzero polynomial with at most $B-1$ monomials
vanishing at $B-1$ positive points. A polynomial with $k$ monomials
has at most $k-1$ distinct positive zeros: factor out its smallest
power and apply Rolle's theorem inductively to its derivative.
This proves the required nonsingularity. For $B=1$, use the empty
determinant $1$ and $w=(1)$.

Form the single sequence
\[
 H_n=\sum_{b=1}^B w_bF_{n,b},\qquad
 q_w(n)=\sum_{b=1}^B w_bq_b(n).
\]
The weights are fixed, so the chosen beta coefficients cancel at
every $n$. All remaining beta values are rational under our
assumption. Let $D_0$ be a common denominator for those rational
values. Proposition~\ref{prop:denominators} and
\eqref{eq:cleared-decay} show that
\begin{equation}\label{eq:integer-sequence}
 D_0\delta_n^S H_n\in\Z,\qquad
 D_0\delta_n^S H_n\longrightarrow0
          \quad(n\to\infty,\ n\text{ odd}).
\end{equation}

The seed vector was chosen to make this sequence nonzero infinitely
often. By \eqref{eq:target},
$q_w(1)=w_B\eta\ne0$. Exclude the finitely many primes dividing
$D_0$ or the numerator or denominator of $q_w(1)$, and take $p$
large enough for Proposition~\ref{prop:dilation}. Its constant
congruence at $m=1$ gives
\[
 p^\tau q_w(p)\equiv\ch(p)2^{-S}q_w(1)\not\equiv0\pmod p.
\]
Every remaining beta term in $p^\tau H_p$ is divisible by $p$ in
$\Zp$, by \eqref{eq:rho-dilation} and the rationality assumption.
Thus $H_p\ne0$ for infinitely many primes $p$.

Equation~\eqref{eq:integer-sequence} now gives infinitely many
nonzero integers tending to zero. Since a nonzero integer has
absolute value at least $1$, this is impossible. Therefore
$I(s)\ge B$, as claimed.
